\documentclass[11pt,a4paper]{article}

% ==============================================================================
% PREÁMBULO ESTANDARIZADO UNSCH - CIENCIAS DE LA SALUD (ENFERMERÍA)
% Compatible tanto con pdflatex como con xelatex (Unicode)
% ==============================================================================

\usepackage{iftex}

\ifPDFTeX
  \usepackage[utf8]{inputenc}
  \usepackage[T1]{fontenc}
  \usepackage{mathptmx} % Fuente Times 12pt requerida por la normativa
\else
  \usepackage{fontspec}
  \setmainfont{Times New Roman}
\fi

\usepackage[spanish,es-nodecimaldot,es-tabla]{babel}

% Márgenes normativos UNSCH:
% Superior: 3.0 cm | Izquierdo: 3.5 cm | Derecho: 2.5 cm | Inferior: 2.5 cm
\usepackage[top=3.0cm, left=3.5cm, right=2.5cm, bottom=2.5cm, a4paper]{geometry}

% Interlineado reglamentario (1.5 líneas)
\usepackage{setspace}
\onehalfspacing

% Sangría reglamentaria de párrafo
\setlength{\parindent}{1.27cm} % Sangría estilo APA (media pulgada)
\setlength{\parskip}{6pt}      % Separación ligera entre párrafos

% Paquetes matemáticos y símbolos
\usepackage{amsmath,amssymb,amsfonts}

% Tablas y matrices avanzadas
\usepackage{booktabs}
\usepackage{tabularx}
\usepackage{array}
\usepackage{multirow}
\usepackage{longtable}

% Inclusión de imágenes y gráficos
\usepackage{graphicx}
\usepackage{float}
\usepackage{caption}
\captionsetup{font=small, labelfont=bf, justification=centering}

% Encabezados y pie de página
\usepackage{fancyhdr}
\pagestyle{fancy}
\fancyhf{}
\fancyhead[R]{\small\thepage}
\fancyhead[L]{\small\nouppercase{\leftmark}}
\renewcommand{\headrulewidth}{0.5pt}
\renewcommand{\footrulewidth}{0pt}
\setlength{\headheight}{15pt}

% Estilo de páginas iniciales
\fancypagestyle{plain}{
    \fancyhf{}
    \fancyfoot[C]{\small\thepage}
    \renewcommand{\headrulewidth}{0pt}
}

% Formato de Títulos
\usepackage{titlesec}
\titleformat{\section}
    {\normalfont\large\bfseries}{\thesection.}{1em}{}
\titleformat{\subsection}
    {\normalfont\normalsize\bfseries}{\thesubsection.}{1em}{}
\titleformat{\subsubsection}
    {\normalfont\normalsize\itshape}{\thesubsubsection.}{1em}{}

% Control de listas numeradas y con viñetas
\usepackage{enumitem}
\setlist[itemize]{noitemsep, topsep=2pt, leftmargin=1.5cm}
\setlist[enumerate]{noitemsep, topsep=2pt, leftmargin=1.5cm}

% Compatibilidad con Pandoc
\providecommand{\tightlist}{%
  \setlength{\itemsep}{0pt}\setlength{\parskip}{0pt}}

% Colores y Enlaces (Modo Tinta Negra Oficial para Tesis y Protocolos UNSCH)
\usepackage{xcolor}

\usepackage{hyperref}
\hypersetup{
    colorlinks=true,
    linkcolor=black,
    citecolor=black,
    urlcolor=black,
    filecolor=black,
    pdfauthor={Gresly Lucero Pariona Palomino},
    pdftitle={Acceso a los servicios de Salud y Calidad de Atencion percibido por los usuarios},
    pdfsubject={Proyecto de Investigacion en Salud (EN 486) - UNSCH},
    pdfkeywords={Acceso a los servicios de salud, Calidad de atencion, Percepcion del usuario, Satisfaccion, Enfermeria, UNSCH, Ayacucho}
}

% Comillas en español
\usepackage{csquotes}

% Microtipografía
\usepackage{microtype}

% Compatibilidad con bloques de codigo de Pandoc
\ifx\Shaded\undefined
  \newenvironment{Shaded}{\begin{quote}\small\ttfamily}{\end{quote}}
\fi
\ifx\Highlighting\undefined
  \newenvironment{Highlighting}{}{}
\fi


\title{\Large\bfseries VIVENCIAS Y SIGNIFICADOS DE LA CALIDAD DE ATENCIÓN EN SALUD: UN ESTUDIO CUALITATIVO EN AYACUCHO}
\author{
    \textbf{Autor Principal}\thanks{Estudiante de la Escuela Profesional de Enfermería, Universidad Nacional de San Cristóbal de Huamanga. ORCID: 0009-0000-0000-0000} \and
    \textbf{Dr. Docente Asesor}\thanks{Docente Principal, Facultad de Ciencias de la Salud, UNSCH. ORCID: 0000-0000-0000-0000}
}
\date{Ayacucho, 2026}

\begin{document}
\maketitle

\begin{abstract}
\noindent\textbf{Objetivo:} Comprender las experiencias del cuidado humanizado de enfermería en pacientes hospitalizados en la ciudad de Ayacucho. \textbf{Metodología:} Estudio cualitativo con diseño fenomenológico hermenéutico. Se entrevistaron a informantes clave mediante entrevistas en profundidad hasta alcanzar la saturación teórica. Los datos fueron examinados mediante análisis de contenido temático respetando los criterios de rigor ético y científico. \textbf{Resultados:} Emergieron tres categorías esenciales: la comunicación empática como bálsamo emocional, el reconocimiento de la dignidad en momentos de vulnerabilidad, y las barreras socioculturales e institucionales para un cuidado cálido. \textbf{Conclusión:} El cuidado humanizado trasciende la destreza técnica, situándose como un acto intersubjetivo vital para la recuperación de la salud.
\vspace{0.2cm}

\noindent\textbf{Palabras clave:} Atención de enfermería; Investigación cualitativa; Cuidado humanizado; Hospitalización.
\end{abstract}

\vspace{0.5cm}

\section{Introducción}
El cuidado constituye el núcleo ontológico y epistemológico de la profesión de enfermería. En escenarios de alta demanda asistencial, la despersonalización del cuidado representa un riesgo constante. La presente investigación analiza la experiencia intersubjetiva de los usuarios en los centros de salud de Ayacucho.

\section{Metodología}
Investigación cualitativa bajo la tradición fenomenológica. La recolección de información se sustentó en entrevistas dialógicas semiestructuradas, audiograbadas y transcritas fidedignamente. Se salvaguardó el anonimato mediante el empleo de seudónimos y la suscripción previa del consentimiento informado.

\section{Resultados y Discusión}
Los relatos evidencian que el gesto, la mirada atenta y la escucha activa de la enfermera modifican positivamente la vivencia de bienestar del paciente, disminuyendo la angustia hospitalaria.

\section{Conclusiones}
Se concluye que fortalecer la formación humanística e intercultural en las aulas universitarias es indispensable para consolidar una práctica de enfermería comprometida con la dignidad del ser humano.

\section*{Referencias}
\begin{enumerate}[label={[\arabic*]}, leftmargin=1cm]
    \item Watson J. \textit{Nursing: The Philosophy and Science of Caring}. Boulder: University Press of Colorado; 2008.
    \item Polit DF, Beck CT. \textit{Investigación en enfermería}. 10.ª ed. Barcelona: Wolters Kluwer; 2021.
\end{enumerate}

\end{document}
